\section{Comparación de casos y patrones de diseño}

\subsection{Desempeño de PPO / SAC en distintas tareas}

Los patrones observados en varios benchmarks son los siguientes:

\begin{itemize}
    \item \textbf{env de baja dimensión y discretos} (CartPole, Acrobot): PPO es suficiente, la entropy de SAC puede causar una exploration variance alta.
    \item \textbf{env de alta dimensión y continuos} (Humanoid, Walker2D): SAC es más estable; SAC + twin Q networks for target smoothing evita la Q divergence.
    \item \textbf{escenarios de cambio de distribución} (Sim2Real): el Replay buffer + prioritized sampling + policy bootstrap layer son clave, la naturaleza off-policy de SAC resulta más ventajosa.
\end{itemize}

\begin{importantbox}{Enfoque híbrido}
En escenarios de Sim2Real u offline RL, la práctica habitual es converger primero la policy con SAC y luego volver a hacer fine-tune con PPO para cumplir con las safety constraints (PPO ofrece clipping guardrails concretos).
\end{importantbox}

\subsection{Patrones de diseño de ingeniería}

Dos patrones habituales:

\begin{enumerate}
    \item \textbf{Dual Agent Pattern}: la policy de PPO y la policy de SAC se entrenan en paralelo; el training job elige automáticamente, según las metrics, la versión más estable para desplegar.
    \item \textbf{Safety Gate Pattern}: se antepone un rule-based gate (supervisory policy) al deployment path; cuando el gate determina que la KL divergence supera el threshold, se pausa el agent y se hace fallback.
\end{enumerate}

\begin{warningbox}{No basta con mirar solo el reward}
En algunos env el reward está diseñado de forma demasiado simple (por ejemplo, basta con completar la tarea para obtener un reward alto), lo que facilita que el agent haga exploit del reward shape. Es indispensable evaluarlo junto con episode length, entropy y KL divergence.
\end{warningbox}

\subsection{Resumen de la sección}

La comparación de casos mostró la estabilidad de PPO en env de baja dimensión y la ventaja de SAC en env de alta dimensión y en Sim2Real, y recomienda usar los patrones dual agent y safety gate para combinar capacidad y controlabilidad.

